% =====================================================================
% Robustness checks
% Status: placeholder structure. No empirical numbers entered.
% Each robustness exercise is the same headline plot
% (chi_distance_curves) recomputed under one perturbation, so the
% reader can compare visually with the main panel of
% Section~\ref{sec:results-distance}.
% =====================================================================
\section{Robustness Checks}\label{sec:robustness}

This section reports the same seasonal dependence--distance
comparison under perturbations of the main methodology. Each
subsection corresponds to one perturbation; all use the same
output format as Section~\ref{sec:results-distance} so that the
reader can compare visually with the main panel.

\subsection{Station selection}\label{sec:rob-stations}

\paragraph{Perturbation.}
Replace the balanced station panel with the unbalanced panel of
Section~\ref{sec:emp-stations}. Pairwise estimators are computed on
pairwise-complete hours, so a station that enters or leaves during
the sample contributes only the pair-hours where both sides are
observed.

\paragraph{What is reported.}
The same dependence--distance figure as
Section~\ref{sec:results-distance}, recomputed on the unbalanced
panel. Also reported: a coastal-vs-inland split, and a sensitivity
panel that varies the within-season completeness threshold of
Section~\ref{sec:emp-missing}.

% TODO: confirm that the unbalanced panel includes at least one
%       station that the balanced panel excludes, so the
%       robustness check has bite.

\paragraph{Interpretation template.}
The text will report whether the seasonal contrast survives the
change in station panel and how much the dependence-distance
curves shift under the relaxed completeness threshold.

\subsection{Extreme-event construction}\label{sec:rob-extremes}

\paragraph{Perturbation.}
Three exercises:
(i) switch the block size in
Section~\ref{sec:meth-extremes} (monthly-within-season
\(\leftrightarrow\) seasonal);
(ii) replace block maxima by threshold exceedances using the POT
route of Section~\ref{sec:meth-extremes-pot}, on a grid of threshold
quantiles;
(iii) under the POT route, vary the minimum cluster separation in
the declustering step.

\paragraph{What is reported.}
A row of dependence--distance panels, one per perturbation, on the
same axes as Section~\ref{sec:results-distance}.

\paragraph{Interpretation template.}
The text will report whether the seasonal contrast is stable
across the perturbations of (i)--(iii), and whether the threshold-
exceedance route gives a qualitatively similar contrast to the
block-maxima route.

\subsection{Marginal standardisation}\label{sec:rob-marginals}

\paragraph{Perturbation.}
Replace the empirical rank transform by the parametric
GEV/GPD tail transform of Section~\ref{sec:meth-marginals}, fit
within each (station, season).

\paragraph{What is reported.}
The headline dependence--distance figure recomputed with the
parametric marginal transform, on the same axes as
Section~\ref{sec:results-distance}.

\paragraph{Interpretation template.}
The text will compare the parametric and rank-transform versions
of the seasonal contrast and flag any pair (or station) where the
two transforms produce qualitatively different dependence
estimates.

\subsection{Sample-splitting in calendar
time}\label{sec:rob-period}

\paragraph{Perturbation.}
Split the sample window into an early and a late half (cut point
to be set in Section~\ref{sec:emp-sample}). Re-estimate
\(\widehat\chi^{(c)}(d)\) on each half separately.

\paragraph{What is reported.}
Two dependence--distance panels (winter and summer), each with
two curves (early, late) and bootstrap bands.

\paragraph{Interpretation template.}
The text will report whether the seasonal contrast is stable
across the two halves of the sample window. This subsection is
the secondary, exploratory long-run analysis announced in
Section~\ref{sec:intro-rq}; the changing station network means
that any observed shift is not directly attributable to climate.

\subsection{Isotropy and direction}\label{sec:rob-direction}

\paragraph{Perturbation.}
Partition the station pairs by the bearing of \(s_j\) relative to
\(s_i\) into a small number of direction bins (e.g.\ four \(45^\circ\)
sectors). Estimate \(\widehat\chi^{(c)}(d)\) within each direction
bin separately.

\paragraph{What is reported.}
A small multiple of dependence--distance panels indexed by direction
bin, one row per season.

\paragraph{Interpretation template.}
The text will assess whether the assumption of isotropy used in
Section~\ref{sec:meth-distance} is innocuous, and whether
prevailing wind direction has a detectable effect on the
dependence--distance summary. Anisotropy is a plausible second-
order effect; we do not model it directly.
